
\documentclass{article}
\usepackage{colortbl}
\usepackage{makecell}
\usepackage{multirow}
\usepackage{supertabular}

\begin{document}

\newcounter{utterance}

\centering \large Interaction Transcript for game `cladder', experiment `full\_v1.5\_default', episode 4371 with reflexion{-}llama{-}3.1{-}8b{-}instruct.
\vspace{24pt}

{ \footnotesize  \setcounter{utterance}{1}
\setlength{\tabcolsep}{0pt}
\begin{supertabular}{c@{$\;$}|p{.15\linewidth}@{}p{.15\linewidth}p{.15\linewidth}p{.15\linewidth}p{.15\linewidth}p{.15\linewidth}}
    \# & \multicolumn{2}{c}{Player} && \multicolumn{2}{c}{Game Master} \\
    \hline

    \theutterance \stepcounter{utterance}  
    & & & \multicolumn{4}{p{0.6\linewidth}}{
        \cellcolor[rgb]{0.9,0.9,0.9}{
            \makecell[{{p{\linewidth}}}]{
                \texttt{\tiny{[P1$\langle$GM]}}
                \texttt{You have attempted this same task before and did not succeed. These are the plans you wrote after those attempts.} \\
\texttt{Trial \#1: New plan: } \\
\\ 
\texttt{In my previous attempt, I tried to estimate the effect of smoking on the probability of being allergic to peanuts using the potential outcomes framework and the counterfactual approach. However, I got stuck in a complex calculation involving multiple conditional probabilities. What I should have done instead is to focus on the direct relationship between smoking and peanut allergy, which is mediated by tar deposit.} \\
\\ 
\texttt{To simplify the analysis, I should have first estimated the effect of smoking on tar deposit, and then used this result to estimate the effect of smoking on peanut allergy. This would have allowed me to avoid the complex calculation involving multiple conditional probabilities.} \\
\\ 
\texttt{Specifically, I should have used the following steps:} \\
\\ 
\texttt{1. Estimate the effect of smoking on tar deposit using the given probabilities: P(T=1|S=1) = 0.58 and P(T=1|S=0) = 0.28.} \\
\texttt{2. Use the estimated effect of smoking on tar deposit to estimate the effect of smoking on peanut allergy, taking into account the direct relationship between tar deposit and peanut allergy.} \\
\\ 
\texttt{By following this simplified approach, I should be able to provide a more accurate and straightforward answer to the question.} \\
\\ 
\\ 
\texttt{You are an expert in causal inference. The following question is not a typical commonsense query, but rather a meticulously designed question created by a professor specializing in causal inference, intended to assess the students' mastery of the course content.} \\
\texttt{Imagine a self{-}contained, hypothetical world with only the following conditions, and without any unmentioned factors or causal relationships: Gender has a direct effect on smoking and peanut allergy. Smoking has a direct effect on tar deposit. Tar deposit has a direct effect on peanut allergy. Gender is unobserved. For nonsmokers, the probability of high tar deposit is 28\%. For smokers, the probability of high tar deposit is 58\%. For nonsmokers and with no tar deposit, the probability of being allergic to peanuts is 49\%. For nonsmokers and with high tar deposit, the probability of being allergic to peanuts is 77\%. For smokers and with no tar deposit, the probability of being allergic to peanuts is 27\%. For smokers and with high tar deposit, the probability of being allergic to peanuts is 62\%. The overall probability of smoking is 58\%. Will smoking increase the chance of being allergic to peanuts?} \\
\texttt{Start your answer with 'Yes' or 'No', followed by additional reasoning or evidence to support your explanation.} \\
            }
        }
    }
     \\ \\

    \theutterance \stepcounter{utterance}  
    & \multicolumn{4}{p{0.6\linewidth}}{
        \cellcolor[rgb]{0.561,0.737,0.561}{
            \makecell[{{p{\linewidth}}}]{
                \texttt{\tiny{[P1$\rangle$GM]}}
                \texttt{Yes, smoking will increase the chance of being allergic to peanuts.} \\
\\ 
\texttt{To answer this question, we need to follow the steps outlined in the plan:} \\
\\ 
\texttt{1. Estimate the effect of smoking on tar deposit:} \\
\\ 
\texttt{Given that P(T=1|S=1) = 0.58 and P(T=1|S=0) = 0.28, we can calculate the effect of smoking on tar deposit as follows:} \\
\\ 
\texttt{P(T=1|S=1) {-} P(T=1|S=0) = 0.58 {-} 0.28 = 0.30} \\
\\ 
\texttt{This means that smoking increases the probability of high tar deposit by 30\%.} \\
\\ 
\texttt{2. Use the estimated effect of smoking on tar deposit to estimate the effect of smoking on peanut allergy:} \\
\\ 
\texttt{We are given the following probabilities:} \\
\\ 
\texttt{{-} For nonsmokers and with no tar deposit, the probability of being allergic to peanuts is 49\%.} \\
\texttt{{-} For nonsmokers and with high tar deposit, the probability of being allergic to peanuts is 77\%.} \\
\texttt{{-} For smokers and with no tar deposit, the probability of being allergic to peanuts is 27\%.} \\
\texttt{{-} For smokers and with high tar deposit, the probability of being allergic to peanuts is 62\%.} \\
\\ 
\texttt{We can use these probabilities to estimate the effect of smoking on peanut allergy as follows:} \\
\\ 
\texttt{Let's define the following variables:} \\
\\ 
\texttt{{-} A: being allergic to peanuts} \\
\texttt{{-} T: high tar deposit} \\
\texttt{{-} S: smoking} \\
\\ 
\texttt{We can use the law of total probability to estimate the overall probability of being allergic to peanuts for smokers and nonsmokers:} \\
\\ 
\texttt{P(A|S=1) = P(A|T=1, S=1) \textbackslash{}* P(T=1|S=1) + P(A|T=0, S=1) \textbackslash{}* P(T=0|S=1)} \\
\texttt{= 0.62 \textbackslash{}* 0.58 + 0.27 \textbackslash{}* (1 {-} 0.58)} \\
\texttt{= 0.62 \textbackslash{}* 0.58 + 0.27 \textbackslash{}* 0.42} \\
\texttt{= 0.3596 + 0.1134} \\
\texttt{= 0.473} \\
\\ 
\texttt{P(A|S=0) = P(A|T=1, S=0) \textbackslash{}* P(T=1|S=0) + P(A|T=0, S=0) \textbackslash{}* P(T=0|S=0)} \\
\texttt{= 0.77 \textbackslash{}* 0.28 + 0.49 \textbackslash{}* (1 {-} 0.28)} \\
\texttt{= 0.77 \textbackslash{}* 0.28 + 0.49 \textbackslash{}* 0.72} \\
\texttt{= 0.2156 + 0.3528} \\
\texttt{= 0.5684} \\
\\ 
\texttt{Now, we can compare the probabilities of being allergic to peanuts for smokers and nonsmokers:} \\
\\ 
\texttt{P(A|S=1) = 0.473} \\
\texttt{P(A|S=0) =} \\
            }
        }
    }
    & & \\ \\

    \theutterance \stepcounter{utterance}  
    & & & \multicolumn{2}{p{0.3\linewidth}}{
        \cellcolor[rgb]{0.9,0.9,0.9}{
            \makecell[{{p{\linewidth}}}]{
                \texttt{\tiny{[GM$|$GM]}}
                \texttt{Yes} \\
            }
        }
    }
    & & \\ \\

    \theutterance \stepcounter{utterance}  
    & & & \multicolumn{2}{p{0.3\linewidth}}{
        \cellcolor[rgb]{0.9,0.9,0.9}{
            \makecell[{{p{\linewidth}}}]{
                \texttt{\tiny{[GM$|$GM]}}
                \texttt{yes} \\
            }
        }
    }
    & & \\ \\

    \theutterance \stepcounter{utterance}  
    & & & \multicolumn{2}{p{0.3\linewidth}}{
        \cellcolor[rgb]{0.9,0.9,0.9}{
            \makecell[{{p{\linewidth}}}]{
                \texttt{\tiny{[GM$|$GM]}}
                \texttt{game\_result = WIN} \\
            }
        }
    }
    & & \\ \\

\end{supertabular}
}

\end{document}
